\section{検証、証明書、残る分類問題}

\subsection{必要 first-jet 方程式の生成}
第21版の全150二次係数から得た固定支持モデルは変更していない。
今回、各成分の normal 座標に関する一次項から $H,B,R$ を計算し、
\[
 \rank B\le1,\quad \tr(H^2)=7R^2/18,\quad
 \det H=-R^3/36,\quad (H+RI/2)B=0
\]
の全係数を保存した。$B=B^{\mathsf T}$ と $\tr H=-R$ は恒等式として確認した。
全ての方程式は、非零な有理数定数による monic 化と完全な重複の除去だけで整理した。
パラメータによる除算、根基、飽和、既約分解はこの生成段階で用いない。

\begin{center}\small
\begin{tabular}{@{}lrrr@{}}\toprule
固定支持モデル & 成分数 & 線形モデルの次元 & 必要 E8 係数式数\\\midrule
$1111$ 鎖 & 4&42&174\\
$1111$ 四重点 &4&48&124\\
$1111$ 三重点と枝 &4&44&157\\
$211$ conic 中央 &3&41&178\\
$211$ 直線中央 &3&41&176\\
$211$ 三重点 &3&43&159\\
$22$ 一節点 &2&40&178\\
$31$ 外部直線 &2&40&186\\\bottomrule
\end{tabular}\end{center}
式数は独立な条件の数、余次元、最小生成元数ではない。
これらの必要式だけで横断型や正規性が十分に判定されるとも主張しない。
実際、前節の非正規な二族も全ての必要式を満たす。

\subsection{異なる実装による厳密な検査}
生成側とは別に、斉次有理座標関数の Jacobian $J$ による $J P J^{\mathsf T}$ を展開し、
normal displacement の一次係数を取って全ての行列を再構成した。
8モデル・25成分の850項目が一致する。

主要な係数証明書は、Python 標準ライブラリだけの別実装でも検証した。
有理数には \texttt{Fraction}、多項式には疎な係数辞書、
有理関数には一変数 Euclid 演算を用いる。
数値代入、浮動小数点、標本点での確認は用いていない。

\begin{center}\small
\begin{tabular}{@{}p{40mm}rp{79mm}@{}}\toprule
独立検証 & 項目数 & 内容\\\midrule
$22$ &228& 元の $P$、局所行列、Hessian、7関係式、4証明書\\
$211$ 三重点 &110& 全3成分の jet と3段の単位イデアル証明書\\
$31$ SS &1243& 七場合、入力式との一致、全消去、残存族の因数分解\\
$31$ NN &658& 全場合、分母因子、二次体、12分岐関係、平面特異族\\\bottomrule
\end{tabular}\end{center}
この標準ライブラリによる検証は計2239項目である。
さらに短い $211$ 証明書36項目、NN 例外族の直接検算35項目、
四重点の錐への帰着12項目、二直線の有限固有値・Jacobi 検算44項目、
局所 weight の検算13項目も行った。
850項目の全座標変換検証と合わせて、追加の symbolic checks は990項目となる。
幾何学的な場合分けの網羅性は本文と独立監査ノートで証明しており、
これらの恒等式チェックがそれに代わるとは主張しない。

\subsection{再現方法と収録物}
\texttt{code/run\_checks.py} は保存済みの係数証明書を厳密に再検査し、
元の模型や前版を書き換えない。出力先は別ディレクトリに指定する。
標準ライブラリの四検証群だけなら Python~3 以外は不要である。
\begin{verbatim}
python3 code/run_checks.py --output-directory /tmp/p4-v22-checks
\end{verbatim}
SymPy~1.14 を用いた全座標変換などの検証も含めるには次を実行する。
\begin{verbatim}
python3 code/run_checks.py --output-directory /tmp/p4-v22-all \
    --symbolic
\end{verbatim}
\path{code/build_e8_jets.py} は8モデルの必要 first-jet 方程式を再生成する。
\path{code/solve_211_triple.py} は三段の証明書を再生成する。
その他の大きな証明書については、全ての入力式、線形結合係数、代入写像を保存し、
標準ライブラリの検証器が恒等式と全特殊値の処理を照合する。

二つの非正規例外族の全 $P$ と $F$ は \texttt{data/families/} に収録した。
本文とは別の幾何学的証明と監査を \texttt{proofs/independent/} に収録した。
旧い探索段階の未成立な推論は今回の証明の入力に含めない。
\texttt{previous\_v21/} 内の全1666ファイルおよび前版 PDF は、保存時の SHA256 と一致する。
\texttt{code/verify\_bundle.py} はこの保存条件と全体の manifest を検証する。

\subsection{分類研究の次の対象}
今回の分類結果は、正規な実際の退化因子を持つ開条件の範囲についての複素閉点の結論である。
全射影 Jacobi スキームについては、次を残す。
\begin{enumerate}
\item 一般階数二：$\pi\wedge\pi=0$ であり、本稿の五次退化因子による議論は適用されない。
\item 一般階数四で、退化因子が被約だが非正規な場合。
可約因子だけでなく、既約で非正規な五次因子も含む。
既知の Feigin--Odesskii の $q_{5,2}$ も非正規な退化因子を持つ\cite{Pym}。
\item 一般階数四で、実際の退化因子が非被約な場合。
\item これらの族の閉包、既約成分、交わり、非被約スキーム構造の決定。
\end{enumerate}
\eqref{eq:ss-factor-family}の二次因子と三次因子を持つ族は、第二の問題へ移るための
明示的な入力になる。一方、この一変数表示が商モジュライの次元であるとは限らず、
その $\PGL_5$ 安定化群と変形空間を調べる必要がある。
\eqref{eq:nn-singular-plane}の二変数族も同じ注意を要する。
今回の八つの排除を繰り返す必要はなく、以後はこれらの残る範囲と
具体的な非正規族の変形・閉包へ進む。
